\section{Actual local growth, admission, and the canonical mean}
\label{sec:local-growth-suppliers-v047}

\noindent\textbf{Scope.}
This chapter supplies the local-growth and mean hypotheses that were
retained in GLG18, and assembles the complete admitted family used
there.  The actual common comparison record
\(\mathrm{GroupCompatibilityFamily}(\mathcal F)\) and the actual
selected-top predicate \(\mathrm{TopsF}(\mathcal F)\) remain explicit
inputs.  The conclusion concerns the same supplied
\(\mathcal F:\mathrm{FamP}(\epsilon)\), its lattices, degrees, support,
and local integer lengths.  We do not construct that common record.

The auxiliary chain introduced below is a variable in the local
counting argument.  It is not a replacement for the geometric theta
chain in the original arithmetic datum.  We verify precisely which
records and conclusions are independent of this auxiliary variable.

\paragraph{Notation and coefficient choices.}
Write \(m=D.R.m\) for the global chain length, \(k=D.R.k\) for the
coefficient rank, \(f=D.R.f\), \(n=D.R.n\), and \(n_v'\) for a local
Frobenius period.  The polynomial rank parameter is \(k-1\), not \(m\).
The source scales are
\[
 \begin{gathered}
 0<\epsilon<1/52,\qquad
 a=xe^{x^\epsilon},\quad X'=e^{x^{\epsilon/3}},
 \quad b=e^{x^\epsilon},\quad b'=a/X',\\
 \delta_1=x^{-\epsilon/8},\quad h=2+8\epsilon,\quad
 r=\lceil8(h+2)/\epsilon\rceil_++1,\\
 G'=\mathrm{Gp},\quad N_1=\mathrm{N1},\quad
 \Lambda=K.\mathrm{La}\,N_1^2\log a,\quad
 p_0=k^2x^{26},\quad c=4\mathrm{Sc}.
 \end{gathered}                                                   \tag{LGS.1}
\]
All ceilings with subscript \(+\) are natural-valued ceilings.
One sufficient original coefficient region is
\[
 \begin{gathered}
 K.\mathrm{La},K.\mathrm{kII},K.\mathrm{tII},
 K.\mathrm{DII},K.\mathrm{szII},K.\mathrm{kII}'\ge4129,\\
 K.\mathrm{kI},K.\mathrm{DI}\ge2800001,\qquad
 K.\rho\ge5600002,\\
 K.\kappa_1\ge604228,\qquad K.\kappa_2\ge302114,
 \qquad K.\mathrm{CN}\ge C_{\mathrm{orb}}(\epsilon).
 \end{gathered}                                                   \tag{LGS.2}
\]
Here \(\kappa_1,\kappa_2\) denote the two original
\texttt{kappaThreshold} coefficients and
\[
 C_{\mathrm{orb}}(\epsilon)
   =40(6r+1)\bigl((h+3)10^{6r+1}+6\bigr)>0.             \tag{LGS.3}
\]
The remaining coefficients are nonnegative.  In the original
\(\mathrm{FamP}\), the fixed fields are
\(K.\mathrm{Pr1}=1\), \(K.\mathrm{dl1}=5\),
\(K.\mathrm{dl2}=4\), and \(\mathcal F.cX.C_e=1\).
We retain them.  The constants in (LGS.2) are chosen before the
variable arithmetic data.  Thresholds in \(x\) follow the fixed
family and all its coefficients, and precede every \(i,u,w,p\),
orbit, coefficient extension, and degree assignment.

% BEGIN LGS1
\paragraph{LGS1. The integral polynomial sandwich at the actual places.}
\label{claim:LGS1}
For every original \(\mathrm{HidaSetting}\), every dominant integral
weight \(\lambda=(\lambda_1,\ldots,\lambda_k)\), and every coefficient
lattice \(A\mathcal O^k\), put
\[
 \begin{split}
 h_\lambda&=\lambda_1-\lambda_k,\qquad
 d_p(\lambda)=\binom{k}{2}\lfloor\log_p(h_\lambda+k-1)\rfloor,\\
 H_\lambda(A\mathcal O^k)
   &=\det(A)^{-\lambda_k}R(A^{-1})
                 \nabla_{\lambda-\lambda_k}(\mathcal O),\\
 D_\lambda(A\mathcal O^k)
   &=\det(A)^{-\lambda_k}R(A^{-1})
                 \Delta_{\lambda-\lambda_k}(\mathcal O).
 \end{split}                                                     \tag{LGS.4}
\]
For rank one the exponent is zero.  If \(h_\lambda=0\), the
inclusion below already holds with exponent zero; the displayed
nonnegative original exponent also works.  Then
\[
       p^{d_p(\lambda)}H_\lambda(A\mathcal O^k)
       \ \subseteq\ D_\lambda(A\mathcal O^k)
       \ \subseteq\ H_\lambda(A\mathcal O^k).            \tag{LGS.5}
\]
This supplies the literal \(\mathrm{ModelSandwichG}\) and
\(\mathrm{WeylSandwichG}\) on the original coefficient DVRs.

Here is the complete matching to the previously proved integral
comparison.  A Hida setting makes \(p\) a nonzero nonunit in
\(\mathcal O\).  Its embedded unramified DVR has characteristic zero:
in positive characteristic \(p\) would be either zero or a unit.
Thus the coefficient fraction field has characteristic zero.
The integer in the source is exactly
\(\mathrm{wExp}(p,k-1,\lambda)
 =N_0(k)\mathrm{deF}(p,k,h_\lambda)=d_p(\lambda)\).
The highest polynomial has coefficient one at its diagonal
monomial.  IW3 therefore applies to the actual coefficient
\(\nabla\) and the actual divided-power span \(\Delta\), with no
rescaling of either lattice or of their highest vector.
It proves \(p^{d_p(\lambda)}\nabla\subseteq\Delta\subseteq\nabla\).

The common invertible fraction-field map in (LGS.4) commutes with
the scalar \(p^{d_p(\lambda)}\), so it transports these inclusions.
Its determinant factor can have negative valuation; it occurs on
both sides and cancels.  Integral changes of basis preserve both
lattices, by CB2 and CR1.  Consequently the same inclusions hold
in the definitions made by intersection over all basis matrices.
For an input that has no presenting lattice basis, the
intersection in each model definition is over an empty
index set, so both are the ambient module and the required
inclusions are automatic.  The separate equivariant-function
definition requires zero values on nonlattice tuples.
For a genuine lattice and a nondominant weight, both
coefficient models are zero.  Thus none of these branches
requires an unverified basis choice.

For the odd Frobenius positions the weight is
\(-\operatorname{rev}\lambda\).  Its spread is again
\(h_\lambda\); (LGS.5) applies with exactly the original exponent.
Tensor products use disjoint polynomial-variable blocks.
Multiplication of polynomials identifies the tensor products
injectively, as may be checked on independent polynomial bases.
Thus exponents add over all positions, without any denominator
from the tensor realization.
% END LGS1

% BEGIN LGS2
\paragraph{LGS2. The original exponent really pays for every factor.}
\label{claim:LGS2}
Let the actual supplied records be
\[
       P:\mathrm{ParameterChoice}'(D,L),\qquad
       T:\mathrm{GatTies}(D).
\]
Suppose
\[
 e_p=X.\mathrm{ep}(p)
       =\left\lceil C_e f k^2\log_p a\right\rceil_+,
       \qquad C_e\ge1.                                  \tag{LGS.6}
\]
Then both fields of the original \(\mathrm{adm9}(D,L,X,C_e)\)
hold:
\[
 \begin{split}
 p\text{ prime},\ p\le a
   &\ \Longrightarrow\
       e_p\le(C_e+1)fk^2\log_p a,\\
 p^{e_p}G.M_{\mathrm{na}}
   &\ \subseteq G.M\subseteq G.M_{\mathrm{na}}
       \quad\text{at every original }(u,w,\mathrm{GIdx}).
 \end{split}                                                     \tag{LGS.7}
\]
The second line holds even when \(p>a\).

We first bound the weight spread on these very factors.
At a selected top the exact endpoint formulas in \(L\) give
\(\lambda_1-\lambda_k=a_t+b_t\).  The top window and
\(k\ge5,\ x^{13}\ge1\) give \(a_t\le6a/5\).
The first nonnegative summand of \(P.\rho_{\mathrm{le}}\)
and \(x^{-2}\le1\) give
\[
 \frac{N_*x^\epsilon\Lambda}{k}\le\frac14,\qquad
 b_t\le\frac54 b'\le\frac54a.
 \]
Hence the top spread is at most \(49a/20<3a\).
Away from tops, the original cloud endpoints give
\(\lambda_1-\lambda_k\le3\mathrm{Sc}\): subtract the lower
last coordinate from the upper first coordinate in the
\(\mathrm{cloud\_place\_bound}\),
\(\mathrm{uk\_extreme}\), and \(\mathrm{cloud\_v}\) formulas;
the \(\rho\)-shift contributes \(-(k-1)\).
The parameter inequality \(k\mathrm{Sc}\le w''\le a\)
therefore also gives a bound \(3a\).
The \(\mathrm{gat\_lam}\) and \(\mathrm{gat\_k}\) identities
transport these statements to every actual \(v,s\).
Odd reversal preserves the spread.

In particular \(0\le h_{\lambda_{v,s}}\le3a\).
The same parameter record gives \(a\ge4\), \(k\le a\), so
\[
       h_{\lambda_{v,s}}+k-1\le4a-1\le a^2.
 \]
Monotonicity of the logarithm and
\(\binom{k}{2}\le k^2/2\) imply
\[
 d_p(\lambda_{v,s})\le k^2\log_p a.                   \tag{LGS.8}
\]
The actual position bijection in \(T.\mathrm{gat\_pos}\) gives
\(\sum_v n_v'=f\), not an inequality with an unspecified
multiplicity.  Summing (LGS.8) yields
\[
 \sum_{v,s}d_p(\lambda_{v,s})
       \le fk^2\log_p a\le C_e fk^2\log_p a\le e_p.
 \]
LGS1 tensorized over precisely these positions now gives the
second line of (LGS.7) at every lattice tuple.  Taking equivariant
functions preserves the inclusions.  The unused exponent is an
additional nonnegative integral power of \(p\).

For \(p\le a\), set \(B=fk^2\log_p a\).  Since \(f\ge1\),
\(k\ge5\), and \(\log_p a\ge1\), one has \(B\ge1\).
The ceiling in (LGS.6) is at most \(C_eB+1\le(C_e+1)B\).
This proves the first line with its original coefficient.
In a \(\mathrm{FamP}\), \(\mathrm{ex.ep\_def}\) gives (LGS.6)
with \(C_e=1\) after the family start.  Thus this is an actual
supplier of \(\mathrm{Adm9F}\) and of
\(\mathrm{LatticeSandwich}\); neither is assumed in the
family application below.
% END LGS2

% BEGIN LGS3
\paragraph{LGS3. Uniform orbit, window, and logarithmic suppliers.}
\label{claim:LGS3}
Every supplied \(\mathcal F:\mathrm{FamP}(\epsilon)\) has the
following original outputs, with \(C_{\mathrm{orb}}\) as in
(LGS.3):
\[
 \begin{gathered}
 \mathrm{Adm1F}\text{ with }C_{31}=C_{\mathrm{orb}}>0,\\
 K.\mathrm{CN}\ge C_{\mathrm{orb}}
       \ \Longrightarrow\ \mathrm{FamilyWindowLogBound},\\
 K.\mathrm{La}>0
       \ \Longrightarrow\ \mathrm{Adm5F}.
 \end{gathered}                                                   \tag{LGS.9}
\]
The first assertion concerns every good orbit of the actual
chosen subset, not a new probabilistic choice.

To check the literal fields, GII1--GII3 give, for \(x\ge3\),
\[
 G'\delta_1\le C_{\mathrm{orb}}\log x,\quad
 N_1\le C_{\mathrm{orb}}\log x,\quad
 \frac{n'\delta_1}{C_{\mathrm{orb}}\log x}
       \le N_v\le40n'\delta_1.                         \tag{LGS.10}
\]
For each selected even position the least positive selected
even successor is within \(G'\) positions and has at most
\(N_1\) interior selected positions.  The number of selected
even predecessors in the preceding \(L_1\) positions is less
than \(r\), because \(L_1\le L_2\) and this is the actual
short-window condition in \(\mathrm{Good}\).
GJ3 identifies every one of these positions, successors and
predecessors with the stored orbit representative.
These are all four conjuncts of \(\mathrm{adm1}\).
They do not require a bound on the proportion of good orbits.
The separate predicate \(\mathrm{adm2}\) will still use the
actual selected-top input.

The middle line of (LGS.9) follows from the second inequality
in (LGS.10).  The definition of \(\rho_B\), nonnegativity of
the coefficients and the actual field
\(\mathcal F.\rho_{\mathrm{le}}\) consequently give
\[
              0\le\rho\le\rho_B\le x^{-2}.             \tag{LGS.11}
\]

For the last line, fix any \(\eta>0\) only after
\(K.\mathrm{La}\) and \(C_{\mathrm{orb}}\).  GJ2 proves all
four original logarithmic inequalities:
\[
 \begin{gathered}
 x^{\epsilon-\eta}\le\Lambda\le x^{\epsilon+\eta},\\
 p<p_0\ \Longrightarrow\ \log(p^{G'})
                       \le x^{\epsilon/8+\eta},\\
 p<X'\ \Longrightarrow\ \log(p^{G'})
                       \le x^{11\epsilon/24+\eta}.
 \end{gathered}                                                   \tag{LGS.12}
\]
For completeness, their single numerical threshold may be
chosen by requiring
\[
 (2K.\mathrm{La}\,C_{\mathrm{orb}}^2
       +38C_{\mathrm{orb}}+(K.\mathrm{La})^{-1})
                  (\log x)^3\le x^\eta,\qquad x\ge3.
 \]
The source bounds \(k\le x^6\), hence
\(\log p_0\le38\log x\); insert (LGS.10) and
\(\log a=\log x+x^\epsilon\).
The positive constant and logarithm powers are absorbed by
the displayed threshold.  It depends on \(\eta\), as the
literal \(\mathrm{Adm5F}\) permits.  \(C_{31}\) does not.
% END LGS3

% BEGIN LGS4
\paragraph{LGS4. The actual small-order prime contribution.}
\label{claim:LGS4}
Every supplied \(\mathrm{FamP}(\epsilon)\) satisfies the
literal \(\mathrm{Adm3F}\):
\[
 \sum_{\substack{p<X''=xnX'\\p\ \mathrm{prime},\
                      \mathrm{IsResidueClassZero}\\
                 \operatorname{ord}(h_P)<x^{\epsilon/4}}}
                    \frac{\log p}{p}
              \le x^{-h+\epsilon}                      \tag{LGS.13}
\]
eventually, uniformly in the family index.

We must first justify use of the dyadic estimate on this
particular Frobenius.  A residue-class-zero prime is odd and
does not divide the host conductor.  The original host
Frobenius fields therefore supply an arithmetic Frobenius
above it.  On the cyclotomic field of conductor \(2^{e+2}\),
arithmetic Frobenius sends \(\zeta\) to \(\zeta^p\).
The host's recorded Galois identification intertwines its
dyadic projection with the class of this \(p\) modulo
\(\{\pm1\}\).  This holds for every permissible prime and
Frobenius, and hence for the particular element used in the
definition of \(h_P\).  In the original notation it says
\[
 D.R.h(h_P)=
   \operatorname{ofMul}\!\bigl(\iota_H(h_{\mathrm{dyad}}(p))\bigr).
                                                               \tag{LGS.14}
\]
For a homomorphism, the order of the image divides the order
of the original element: apply the homomorphism to its
annihilating multiple.  Thus
\(\operatorname{ord}(h_{\mathrm{dyad}}(p))
 \le\operatorname{ord}(h_P)\).
This supplies the exact \(\mathrm{FrobeniusClassCompatibility}\)
used here; it is not a replacement of \(h_P\) by a chosen
conjugacy representative.

Put \(M=2^e\) and \(J=x^{\epsilon/4}\).
The rank bounds below give \(M\ge x^{h-\epsilon/50}>1\)
once \(x>1\), so the dyadic exponent satisfies \(e\ge1\).
The integer dyadic orders in this sum are at most
\(\lfloor J\rfloor\).  LFS2, with
\(\lfloor J\rfloor\le J\), gives
\[
 \sum\frac{\log p}{p}
       \le\frac{12J\log^2(MX'')}{M}.                  \tag{LGS.15}
\]
The source rank relations and its \(\mathrm{SizeN}\) and
\(\mathrm{o\_tie}\) fields are
\[
 M\le n\le oM,\quad o\le x^{\epsilon/100},\quad
 x^{h-\epsilon/100}\le n\le x^{h+\epsilon/100}.
 \]
Consequently \(M\ge x^{h-\epsilon/50}\) and
\(\log M\le\log n\le3\log x\).
Choose the threshold so that
\[
 8\log x\le x^{\epsilon/3},\qquad
                  48\le x^{19\epsilon/300}.
 \]
Then
\[
 0\le\log(MX'')\le
       7\log x+x^{\epsilon/3}\le2x^{\epsilon/3}.
 \]
The right side of (LGS.15) is at most
\[
 48x^{-h+(1/4+2/3+1/50)\epsilon}
 =48x^{-h+281\epsilon/300}
 \le x^{-h+\epsilon}.
 \]
Both thresholds precede every prime and every member.
The actual sums are finite; using \(p<X''\) instead of
\(p\le X''\) only restricts the summation set.
No new prime-distribution theorem is required.
% END LGS4

% BEGIN LGS5
\paragraph{LGS5. All five original power-band inequalities.}
\label{claim:LGS5}
Under \(K.\mathrm{CN}\ge C_{\mathrm{orb}}\), the actual family
satisfies \(\mathrm{Adm4F}\), including both prime sums with
their original bound \(x^{-1}\).

We prove the finite estimate used for both bands.  Let
\(X\ge1\), \(P_0\ge1\), \(0\le q\le1/3\), \(d\ge1\), and
let \(S\) be any finite set of positive integers such that
\[
 P_0\le p<X,\qquad X(1-q)\le p^d\le X(1+q)
                       \quad(p\in S).
 \]
If \(S\ne\varnothing\), let \(u=\min S\), \(v=\max S\).
Since \(v/u\ge1\),
\[
 \frac vu\le\left(\frac vu\right)^d
       \le\frac{1+q}{1-q}\le1+3q.
 \]
There are at most \(v-u+1\le3qu+1\) integers between the
endpoints, while \(\log p/p\le\log X/u\).
Thus, including the empty case,
\[
             \sum_{p\in S}\frac{\log p}{p}
                 \le\log X(3q+P_0^{-1}).             \tag{LGS.16}
\]
If each \(p\) is prime and some positive exponent \(d\) works,
choose one such exponent per \(p\).  As \(p\ge2\) and
\(p^d\le2X\),
\[
       d\le1+\frac{\log X}{\log2}\le1+2\log X.
 \]
Partition the finite set by the chosen exponent and apply
(LGS.16) on each part.  This yields
\[
 \sum_{p\in S}\frac{\log p}{p}
       \le(1+2\log X)\log X(3q+P_0^{-1}).             \tag{LGS.17}
\]
This proof counts an endpoint integer once even for a
zero-width band; it does not discard the \(P_0^{-1}\) term.

Apply (LGS.17) with \(X=X'\), \(P_0=p_0\), and first
\(q=\rho\), then \(q=\rho_B\).  LGS3 gives
\(0\le q\le x^{-2}\).  For \(x\ge2\) this is at most \(1/4\),
and \(p_0=k^2x^{26}\ge x^2\).  With
\(B=\log X'=x^{\epsilon/3}\ge1\), the right side is at most
\[
       (3B)B(4x^{-2})=12x^{2\epsilon/3-2}\le x^{-1}
                                                               \tag{LGS.18}
\]
as soon as \(12\le x^{1-2\epsilon/3}\).
The exponent is positive since \(\epsilon<1/52\).
Each original \(\mathrm{InBand0}\) or \(\mathrm{InBand}\)
prime satisfies precisely the premises of this finite
estimate.  Together with (LGS.11), the two bounds in
(LGS.18) are all five fields of the original
\(\mathrm{ArchimedeanSignatureMatching}\).
They use one common threshold, independent of the prime
and its witnessing exponent.
% END LGS5

% BEGIN LGS6
\paragraph{LGS6. The actual two projection certificates and canonical mean.}
\label{claim:LGS6}
Under (LGS.2), the same instance \(X=\mathcal F.Gx(x,i)\)
eventually supplies all three original predicates
\[
 \mathrm{adm6}(D,L,K),\qquad
 \mathrm{adm7}(D,L,X,K),\qquad
 \mathrm{adm8}(D,L,X,K,302114).                         \tag{LGS.19}
\]
In addition, the actual \(\mathrm{HypIAt}\) and
\(\mathrm{HypIIAt}\) hold at every permitted prime and good
orbit.  All coefficient settings and dominant weights in
these predicates retain their original universal quantifiers.

Here are the numerical and model identifications needed to use
the completed local theorems.  LGS3 supplies the same
\(C_{31}\), window count and \(\mathrm{adm5}\) with
\(\eta=\epsilon/48\).  LGS5 supplies exclusion of the original
power band.  In the large-prime range the fourteen inequalities
of the source's \(\mathrm{HypI}\), together with the two
spacing inequalities and actual degree windows, are exactly
GI1--GI5.  Their coefficients are
\[
 C_H=2800001,\quad C=302114,\quad C_0=K.\mathrm{La},\quad
 m_{\mathrm{loc}}=k-1,\quad c_{\mathrm{loc}}=\lfloor4\mathrm{Sc}\rfloor_+.
                                                               \tag{LGS.20}
\]
The lower bounds \(K.\mathrm{kI},K.\mathrm{DI}\ge C_H\),
\(K.\rho\ge2C_H\), \(K.\kappa_1\ge2C\),
\(K.\kappa_2\ge C\) and \(K.\mathrm{La}\ge880\)
are all consequences of (LGS.2).  The candidate excludes the
band of radius \(\rho_B\); LGS3 gives \(\rho\le\rho_B\),
so it also excludes the local forbidden band.  The source
windows are the literal restrictions of \(a_L,b_L\) to
the cyclic orbit, by GJ3.  No recentering is performed.

In the small-prime range GII4--GII7 supply every numerical
and geometric hypothesis of the Type-II theorem on the same
orbit.  GII8 applies with the single original constant \(4129\),
since it is no larger than
\[
 K.\mathrm{kII},\ K.\mathrm{tII},\ K.\mathrm{La},\
 K.\mathrm{DII},\ K.\mathrm{szII},\ K.\mathrm{kII}'.
\]
This is precisely the full \(\mathrm{adm6}\); it is not
restricted to coefficient rings occurring at a chosen
global place.  In particular it remains valid after the
standard-ring construction used below.

For \(\mathrm{adm7}\), the supplied family field
\(\mathrm{ex.kapI\_def}\) says that \(X.\kappa_I\) is the
infimum of the deficits for which the required certificate
holds universally over those coefficient settings.
GJ1 proves attainment.  To recall why this has no compactness
gap, in one fixed setting a certificate of real rank at least
\((1-q)d\) has an integral number of independent tests.
For \(q\) sufficiently close above its infimum \(q_0\),
that integer is at least \(\lceil(1-q_0)d\rceil\).
All other tests and divisibilities are unchanged.  The
setting was arbitrary, so \(q_0\) itself satisfies the
universal predicate.  This proves \(\mathrm{adm7}\) for the
original \(X.\kappa_I\), rather than a newly chosen function.

Finally GI6--GI8 compare the local theorem's time with the
original \(\lceil t_{If}\rceil_+\), retaining the common
signed leading sum.  Its exceptional function is chosen
before the coefficient setting and has mean at most
\[
                   302114\,N_vN_1/x^{13}.             \tag{LGS.21}
\]
Leastness of the attained canonical deficit bounds
\(X.\kappa_I\) by this same function.  Restriction to the
selected orbit positions is a bijection between the
original degree box and the local product box, with inverse
extension by zero.  Thus both sums and both cardinality
denominators agree, proving \(\mathrm{adm8}\) with
\(C_m=302114\).  Empty boxes obey the original zero-average
convention.  Taking the maximum of the thresholds just
described proves (LGS.19) uniformly in \(i\), before any
prime, degree assignment or coefficient extension is chosen.
% END LGS6

% BEGIN LGS7
\paragraph{LGS7. The literal mean bound on the original canonical deficit.}
\label{claim:LGS7}
Under (LGS.2), the supplied family eventually satisfies
\(\mathrm{MeanKappaBound}(\mathcal F.Gx,K)\).
Explicitly, for every candidate prime and every actual good
representative,
\[
 \begin{split}
 \operatorname{meanLoc}(\kappa_{\mathrm{loc}})
       &\le N_vx^{-12},\\
 p<p_0\quad&\Longrightarrow\quad
       \kappa_v(\mathbf a,\mathbf b)\le N_vx^{-12}
       \quad\hbox{for every }\mathbf a,\mathbf b.
 \end{split}                                                     \tag{LGS.22}
\]
This supplies both clauses, including the pointwise
small-prime clause, on the same \(X\) as LGS6.

Indeed the definition in the small-prime range is
\[
 \kappa_{\mathrm{loc}}=\kappa_v
     =\min\{1,K.\mathrm{kII}\,N_vr^2k x^{-26}\}.
 \]
It is independent of the degrees and nonnegative.  Since
\(k\le x^6\), the single requirement
\(K.\mathrm{kII}\,r^2\le x^8\) makes it at most
\(N_vx^{-12}\).  Its average has the same bound; if its
box is empty, the average is zero.
For \(p\ge p_0\), the candidate's order lower bound and
LGS3 supply the orbit, spacing and window hypotheses of
LGS6.  Consequently (LGS.21), together with
\[
                302114\,C_{\mathrm{orb}}\log x\le x,   \tag{LGS.23}
\]
gives the first line of (LGS.22).  Both requirements are
eventual in \(x\) after the fixed coefficients.  They are
independent of \(p,v,i\), of the coefficient DVR, and of
the degrees.  This is the numerical content of GW1--GW2
matched to the full original \(\mathrm{MeanKappaBound}\).

In particular no assumption that the actual selected
degrees already have a small deficit at every prime has
been inserted.  The later finite averaging theorem can
use (LGS.22) to choose a single degree assignment; it does
not change the function proved to satisfy this mean bound.
% END LGS7

% BEGIN LGS8
\paragraph{LGS8. What the common compatibility record actually supplies.}
\label{claim:LGS8}
Let the original \(\mathrm{GroupCompatibilityFamily}(\mathcal F)\)
be supplied.  Then, after one family threshold, the actual
local models have finite class sets and trivial stabilizers,
the standard coefficient-subring interface
\(\mathrm{StdSubAt}\), and \(\mathrm{NPDescF}\).
Both \(G.M_{\mathrm{na}}\) and \(G.M\) are finite free of the
same rank \(G.d\).  The two original operator-stability
conditions hold on these very lattices.

The matching uses every part of the record in its intended
role.  At fixed \(x,i\), its fields are
\[
 \mathrm{herm},\quad \mathrm{forms},\quad
 \mathrm{cmp},\quad \mathrm{ring},\quad \mathrm{gat}.    \tag{LGS.24}
\]
The last field identifies \(D.gat(u,w,h_g)\) with the
actual \(\mathrm{gatOfLevelP}\) and includes the original
\(\mathrm{Cmp.HypAtP}\); the level is the very
\(\mathrm{herm.lev0}(u)\).  It is not merely an isomorphism
of two abstract representation spaces.
The two distinct auxiliary primes and their
upper-unitriangular level conditions are part of this
level datum.  CR3 identifies its lattice orbits with its
finite class set, and CR4 makes the stabilizers trivial.
CR2 then identifies integral equivariant functions with
the finite product of their actual fibers.  CR1 gives
finite free Weyl fibers; CB2 gives finite free
dual-Weyl fibers.  LGS1 makes their rational spans equal.
Their common function-space rank is therefore exactly
\(G.d\), with the original character and determinant
twists.  This proves the literal
\(\mathrm{TrivialStabiliserFiniteClasses}\), not just a
finite-dimensionality assertion over the fraction field.

For the standard subring, \(\mathrm{ring}\) supplies four
compatible isomorphisms \(W,Q,\mathcal O,K\), including the
maps of the unramified subring, its fraction field, and
the chosen square root of \(p\).  Compose the finite map
constructed in BC1--BC2 with these isomorphisms.
BC8 shows that its three compatibility equations are
preserved.  Thus the target is the actual coefficient
DVR and the image of the chosen square root is the
recorded \(w\), even if it is not a uniformizer.
BC6--BC7 consequently transport the local rank
certificates without replacing the residue cardinality
of the unramified ring or changing any \(p\)-normalized
depth.

Finally the \(\mathrm{cmp}\) field is the same
\(\mathrm{CmpTied}\) used with the given
\(\mathrm{GatTies}\) and the family's actual comparison
map.  GLG14 applies its rational retraction and faithful
flatness argument to the ordered products and ramified
intersection, proving \(\mathrm{NPDescF}\).
GLG13 proves both operator-stability statements directly
for the actual polynomial models, including signed
determinant factors and nondominant zero fibers.
No extra integrality or comparison premise remains at
these steps.  Construction of (LGS.24) itself is the
explicit retained input of this chapter.
% END LGS8

% BEGIN LGS9
\paragraph{LGS9. The actual orbit certificates give joint exceptional generators.}
\label{claim:LGS9}
Assume (LGS.2) and the common compatibility record.
At every actual \((u,w,h_g)\), every served representative
\(t_0\) supplies the original \(\mathrm{LPprime}\) with
\[
 \begin{gathered}
 v=vOf(t_0),\qquad t_v=\lceil t_v^{\mathrm{source}}\rceil_+,
 \qquad D_0=D_p,\qquad \kappa_v=\kappa_v^{\mathrm{source}},\\
 c_*=8(D_p+1),\qquad
 r_{\mathrm{loc}}=
       \begin{cases}r,&p<p_0,\\1,&p\ge p_0.\end{cases}
 \end{gathered}                                                   \tag{LGS.25}
\]
It supplies \(\mathrm{OrbitLatticeBounds}\) for the original
data, not just for isomorphic standard models.
For every nonempty subset \(B\) of the served local
coordinates, the actual dual-slot section lattice \(L_B\)
has a finite set \(E_B\subset L_B\) such that
\[
 \begin{split}
 |E_B|&\le d\prod_{v\in B}\kappa_v,\qquad d=G.d,\\
 p^{D_p+1}L_B
   &\subseteq
     \sum_{v\in B}y_v L_{B;\,v=\mathrm{ball}}
                         +\langle E_B\rangle_{\mathcal O}.
 \end{split}                                                     \tag{LGS.26}
\]

For (LGS.25), a served representative belongs to the actual
good representative set of a candidate prime.  Its cyclic
coordinate bijection is GJ3.  The \(\mathrm{LatTied}\)
equalities give the same period, selected positions,
top and bottom degrees, filling weights and local rank
parameter.  The remaining coordinate bound is at most
\(\lfloor4\mathrm{Sc}\rfloor_+\).
LGS6 provides the applicable Type-II or Type-I certificate
over the standard subring; LGS8 and BC6--BC7 transport it
to the actual \(\mathcal O\).  GJ4 checks the four
quantities in (LGS.25) exactly.  Thus no height, depth or
exceptional-rank allowance is enlarged in this transport.

We now check all geometric hypotheses for the product
estimate in (LGS.26).  In each coefficient factor, EL5--EL7
construct the source, midpoint and target lattices from an
actual edge in Smith coordinates.  With
\[
 A_v=\sum_s\sum_{j=1}^{r_{\mathrm{loc}}}\lambda_{s,j},
 \]
the normalized edge pairing is
\(\beta_e(F,\ell)=\ell(p^{A_v}F)\).
Its integral perfection follows from
\(w^{-A_v}p^{A_v}w^{-A_v}=1\); \(A_v\) may be negative.
The ball and dual slots are the actual intersections and
polars in EL7, so for fixed other slots
\[
 P_t\subseteq H\subseteq B_t,\qquad
 w^tB_t\subseteq H,\quad w^tH\subseteq P_t,\quad
 y_v(B_t)\subseteq P_t.                                \tag{LGS.27}
\]
EL8--EL10 prove their equivariance and transport the
unwanted-edge errors.  These statements are valid after
tensoring all unchanged coordinates and after taking
the actual equivariant sections.  All slot lattices are
finite free and full by LGS8 and their two-sided
comparison with the base fibers.
The map \(vOf\) is injective on the representative set
by \(\mathrm{LatTied}\), so distinct representatives
produce distinct coordinates.  Operators at distinct
coordinates commute: both compositions are the same
finite double sum over independent incoming edges.
Their scalar factors also commute.

For clarity, the integral rank argument is as follows.
The \(\mathrm{LPprime}\) tests give an integral matrix
identity equal to \(p^{D_p}I_n\), with
\(n\ge(1-\kappa_v)d_v\), on the local dual-slot lattice.
JR4 applies Smith form and finds at least \(n\) basis
directions divisible by these tests at depth \(D_p\).
At most \(\kappa_vd_v\) directions are bad.
At an actual orbit representative tensor the bases and
retain a tensor only if every selected coordinate is bad.
Its count is at most
\((\prod_{v\in B}\kappa_v)\prod_vd_v\).
The finite orbit and trivial-stabilizer facts of LGS8
extend each retained value uniquely to a delta section.
The total base rank is \(d\), so the count is the first
line of (LGS.26).

Every discarded tensor has a good selected coordinate.
Applying its neighbor relation gives the required value
once at that target, since its stabilizer is trivial.
All unintended targets lie in the stated
\(w^{c_*}\)-error lattice; EL10 supplies the integral
identification with the target dual slot.
Here \(c_*=8(D_p+1)\ge4(D_p+1)\), so the error has
at least the \(p^{D_p+1}\) divisibility used in JR5.
The resulting finite quotient satisfies \(Q=pQ\);
the elementary determinant proof of Nakayama, or JR7,
makes it zero.  This proves the second line of
(LGS.26) with the original exponent \(D_p+1\).
The argument has now supplied the quantitative,
polynomial, orbit, and stabilizer inputs that were
separate in JR1--JR12 and LW4.
% END LGS9

% BEGIN LGS10
\paragraph{LGS10. The literal first-growth inclusion and length.}
\label{claim:LGS10}
Under the same hypotheses, fix any actual place and word.
Put \(M=G.M\), \(M_{\mathrm{na}}=G.M_{\mathrm{na}}\),
\(U=Y.Vram\), and \(M_{\mathrm{ram}}=M\cap U\).
Then eventually, with the exact original \(\tau_P\),
\(\mathrm{io}_P\), and \(\mathrm{NP}\),
\[
 \begin{split}
 \mathrm{NP}&\subseteq p^{\lceil\tau_P\rceil_+}M_{\mathrm{ram}},\\
 \ell:=\operatorname{length}_{\mathcal O}
                 (M_{\mathrm{ram}}/\mathrm{NP})&\in\mathbb N,
 \qquad
 \ell\le vUnit(G)\,G.d\,(\tau_P+\mathrm{io}_P).
 \end{split}                                                     \tag{LGS.28}
\]
These are the first two conjuncts of the original
\(\mathrm{admG1}\), including the finite natural-valued
length, for every \(\mathrm{GatTies}\) and
\(\mathrm{LatTied}\) witness on the fixed family data.

Let \(\mathcal V\) denote the actual served set.
If it is empty, the empty product term in the definition
of \(\mathrm{NP}\) is \(M\), while
\(\tau_P=\mathrm{io}_P=0\).  Hence
\(\mathrm{NP}=M_{\mathrm{ram}}\) and (LGS.28) follows.
Suppose it is nonempty.  Set
\[
 \begin{gathered}
 S_v=\mathrm{Sv}(v),\quad S=\max_{v\in\mathcal V}S_v,\quad
 \Sigma=m\sum_{v\in\mathcal V}\kappa_v,\\
 \tau_0=\sum_{v\in\mathcal V}m(t_v+D_p+1),\qquad
 \tau=e_p+\tau_0 .
 \end{gathered}                                                   \tag{LGS.29}
\]
Each \(S_v\) is already the integer ceiling of a
nonnegative sum, so is a nonnegative integer.
The source's outer natural-valued ceiling therefore
does not add one.  The source definition similarly
gives \(\tau_P=\tau\) exactly.
Because a nonempty served set is precisely the whole
good representative set at a candidate with
\(\mathrm{Sig}\le x^{-\epsilon}\), one has
\[
       \Sigma=\mathrm{Sig}\le x^{-\epsilon}\le\tfrac12,
       \qquad \mathrm{io}_P=2S\Sigma                  \tag{LGS.30}
\]
as soon as \(x\ge2^{1/\epsilon}\).
The original rank record already gives \(m\ge2\).

Use the actual slot lattices and joint relations of
LGS9 in the finite walk LW1--LW3.  For
\[
 J_{\mathrm{na}}=
 \sum_{B\subseteq\mathcal V}
     p^{\sum_{v\notin B}S_v}
       \Bigl(\prod_{v\in B}y_v^m\Bigr)M_{\mathrm{na}},
 \]
it gives, for every integer \(\beta\ge0\), at most
\[
               d\,m^{\beta+1}e_{\beta+1}(\kappa)       \tag{LGS.31}
\]
exceptional directions at depth \(\tau_0+\beta S\).
One sees the multiplicity directly in that walk:
choose the \(\beta+1\) active coordinates and one of
the \(m\) counter positions at each active coordinate.
Its depth cost at a coordinate is exactly
\(m(t_v+D_p+1)\).  Thus neither \(m\) nor the \(+1\)
is omitted.

The actual sandwich LGS2 and operator stability LGS8
give \(p^{e_p}J_{\mathrm{na}}\subseteq J\), where \(J\)
is defined by the same formula with \(M\).
Multiply the walk inclusion by \(p^{e_p}\).
The exceptional vectors then lie in \(M\), their number
is still (LGS.31), and their depth in \(M\) is
\(\tau+\beta S\), exactly as in SR3.
In particular the empty product term shows that \(J\)
contains a nonzero scalar multiple of \(M\), so all
lengths in this calculation are finite.

Write \(e_{\mathcal O}=v_{\mathcal O}(p)=vUnit(G)>0\).
By LC1--LC2, summing the Smith tails gives
\[
 \begin{split}
 \operatorname{length}_{\mathcal O}
       \bigl(M/(J\cap p^\tau M)\bigr)
 &\le e_{\mathcal O}d\left[
       \tau+S\left(\prod_v(1+m\kappa_v)-1\right)\right]\\
 &\le e_{\mathcal O}d(\tau+2S\Sigma).
 \end{split}                                                     \tag{LGS.32}
\]
The last inequality follows from
\(\prod_v(1+m\kappa_v)\le e^\Sigma\) and
\(e^\Sigma-1\le2\Sigma\) for \(0\le\Sigma\le1/2\).
For \(S=0\), the vanishing of (LGS.31) past
\(|\mathcal V|\) already puts every Smith exponent
at depth at most \(\tau\); the same formula holds.

The actual field \(\mathrm{ex.lat.vram}\) identifies
\(U\) with the coefficient-field span of the original
ramified subspace.  If \(a\ne0\) and \(az\in U\), its
inverse in that field gives \(z\in U\).  Therefore
\(M_{\mathrm{ram}}\) is saturated in \(M\), and
\(p^\tau M\cap M_{\mathrm{ram}}=p^\tau M_{\mathrm{ram}}\).
Now
\[
 \mathrm{NP}=J\cap p^\tau M\cap U.
 \]
This proves the inclusion in (LGS.28); the quotient
\(M_{\mathrm{ram}}/\mathrm{NP}\) injects into the
finite quotient in (LGS.32).  Its natural-valued length
therefore obeys the same bound.  Substituting (LGS.29)
and (LGS.30) proves (LGS.28) with the original full
rank \(d=G.d\) and normalization \(vUnit(G)\).
% END LGS10

% BEGIN LGS11
\paragraph{LGS11. An auxiliary zero chain preserves the complete interface.}
\label{claim:LGS11}
Define \(Y^0\) from the original \(\mathrm{LatInst}\)
\(Y\) by keeping \(vOf,Fl,Lv,Vram\) and replacing every
linear map of its chain \(Ch\) by zero.
Then the same family eventually satisfies
\[
 \begin{gathered}
 \mathrm{unitBoundAtTwistedPlaces}(D,L,X,Y^0,K),\\
 \mathrm{UnitBoundDifferenceTerm}(D,L,X,Y^0,K),\qquad
 \mathrm{admG1}(D,L,X,Y^0,K).
 \end{gathered}                                                   \tag{LGS.33}
\]
This construction requires no lower-valuation estimate
for the original theta chain.

We check the record rather than only its final numerical
expression.  A \(\mathrm{Chain}\) consists solely of
coefficient-field linear maps \(K_j\) between the
successive spaces \(Lv.Fm(j)\).  Zero maps satisfy that
definition.  Since \(m\ge2\), their composite
\(K_{m-1}\cdots K_0\) is zero on every input.
The predicate \(Y.vKge(F,g)\) asks for a uniformizer
\(\varpi\) such that this composite lies in
\[
 \varpi^{\lceil g\,vUnit(G)\rceil}\,Lv.Int(m).
 \]
A DVR has a uniformizer, and zero lies in this scaled
module for every real \(g\), including negative \(g\).
Thus \(Y^0.vKge(F,g)\) holds for every \(F,g\).

Every clause of \(\mathrm{LatTied}\) uses only \(vOf\),
\(Fl\), and the fixed \(D,X\).  Its period, rank,
selected-coordinate, filling, and injection conditions
are consequently unchanged.  The ordered product
\(\mathrm{yProd}\) uses only \(vOf\) and the actual
\(G.y_v^m\), so it too is unchanged.  In the complete
definition of \(\mathrm{NP}\), the three factors are
the span of these products on \(M\), the scalar-depth
copy of \(M\), and \(Vram\).  Hence
\[
 \mathrm{NP}(X,Y^0,K)=\mathrm{NP}(X,Y,K),\qquad
 M_{\mathrm{ram}}(Y^0)=M_{\mathrm{ram}}(Y).             \tag{LGS.34}
\]
The rational descent and scalar-extension identities
transfer exactly.  So does every original
\(\mathrm{LatticeLengthsAtPlace}\) identity.  The numerical
functions \(\kappa_v,S_v,\mathrm{Sig}\), \(VV\),
\(\tau_P,\mathrm{io}_P\), and \(\mathrm{gainP}\)
do not mention \(Ch\).
Neither do \(\mathrm{ParametersFill}\),
\(\mathrm{NuDef}\), the support condition, or any of
the five fields of \(\mathrm{LinkG\_at}\).
The actual selected global lattice \(L.NuLat\) and
its comparison with \(M_{\mathrm{ram}}\) are unchanged.
More explicitly, \(\mathrm{LatDefs}\) has just the
\(\mathrm{vram}\), \(\mathrm{served}\), and
\(\mathrm{atT}\) fields: the first two transfer by
(LGS.34), and the third does not mention \(Y\).
The larger arithmetic \(\mathrm{Extra}\) record
also has a \(\mathrm{ker}\) field about its geometric
chain.  That record remains the original record;
we do not transport its \(\mathrm{ker}\) field to
\(Y^0\).  Neither auxiliary admission nor the
\(\mathrm{LinkG}\) installer requires that transport.

The first two lines of (LGS.33) are now immediate even
with their full quantifiers: for every pair of subsets
\(VV'\subseteq VVs\subseteq\mathrm{reps}\), every
original section, and the printed lower term, the
left composite is zero.  The same observation gives
the final valuation conjunct of \(\mathrm{admG1}\)
for every \(F\in\mathrm{NP}\) and the exact lower term
\(\sum_j\mathrm{gainP}_j+\mathrm{lowFE}\).
The first two conjuncts are LGS10, which apply to the
unchanged \(\mathrm{NP}\) and \(M_{\mathrm{ram}}\).
This proves the full original \(\mathrm{admG1}\) on
the auxiliary \(Y^0\).

Only this auxiliary argument has a zero chain.
The given arithmetic datum, its geometric chain,
its valuations, and all divisibility hypotheses used
elsewhere have not been altered.  Nothing in
(LGS.33) asserts a new estimate for those original
theta maps.
% END LGS11

% BEGIN LGS12
\paragraph{LGS12. The complete original admitted family is now supplied.}
\label{claim:LGS12}
Assume (LGS.2), the original
\(\mathrm{GroupCompatibilityFamily}(\mathcal F)\), and
the actual \(\mathrm{TopsF}(\mathcal F)\).
Then
\[
 \mathrm{AuxBlock8.AdmF}
   \bigl(\mathcal F,K,\mathcal F.Gx,
                     (x,i)\mapsto(\mathcal F.Gy(x,i))^0\bigr)      \tag{LGS.35}
\]
holds.  The constants may be fixed as
\[
       C_{31}=C_{\mathrm{orb}}(\epsilon),\qquad
                         C_e=1,\qquad C_m=302114.     \tag{LGS.36}
\]
Equivalently, these same constants supply the full
\(\mathrm{AdmFam}\).  The assumption \(\mathrm{TopsF}\)
may instead be replaced by the literal \(\mathrm{Adm2F}\).

We give every field and its threshold, since a partial
admission record would not suffice here.
LGS3 gives \(\mathrm{adm1}\) with this \(C_{31}>0\).
GLG16 applies to the original selected tops and gives
the full \(\mathrm{adm2}\): both its cardinality
inequalities and its weighted bad-top bound.
It takes a threshold after the one actual constant
in \(\mathrm{ChosenBad}\), without resampling the tops.
LGS4 and LGS5 give \(\mathrm{adm3}\) and all five
fields of \(\mathrm{adm4}\).
For any \(\delta>0\), LGS3 gives \(\mathrm{adm5}\)
with that \(\delta\).  Its threshold may depend on
\(\delta\), and (LGS.36) does not.
LGS6 gives \(\mathrm{adm6},\mathrm{adm7},\mathrm{adm8}\)
with the canonical deficit and the constant \(C_m\)
in (LGS.36).  LGS2 gives \(\mathrm{adm9}\) with
the actual fixed \(C_e=1\).
LGS11 supplies the unit bound and every conjunct of
\(\mathrm{admG1}\) for \(Y^0\).

It remains to check the original
\(\mathrm{TypePrimeSizeConds}\), which has three fields.
The family's \(\mathrm{sizes}\) field gives all of
\(\mathrm{SizesAt}(L)\).
The actual parameter record and \(\mathrm{fills}\)
give, respectively, its \(S\)-bound and the equality
of the two \(\Lambda\)'s, hence
\[
                  4\mathrm{Sc}\le\Lambda b'/k.
 \]
GLG1 gives \(N_m\le\mathrm{Sc}\),
\(\mathrm{Sc}\ge100\), \(k\ge5\) and
\(k\mathrm{Sc}\le w''\).  Therefore
\(N_m+1\le\mathrm{Sc}+1\le k\mathrm{Sc}\le w''\).
This is the third field, with the original \(N_m\).

Fix (LGS.36) first.  Given \(\delta>0\), choose the
maximum of the finitely many preceding thresholds,
the compatibility threshold, the actual family start
plus one, \(2^{1/\epsilon}\), and \(2\).
For every later \(x\) and every \(i\in\mathcal F.Ix(x)\)
all twelve fields hold simultaneously.  This is
exactly
\[
 \exists C_{31}>0,\ C_e,C_m\ge0\
 \ \forall\delta>0\
 \ \exists x_0>1\
 \ \forall x\ge x_0\
 \ \forall i,\quad
             \mathrm{admitted}(D,L,X,Y^0,K;
                                      C_{31},C_e,C_m,\delta).
                                                               \tag{LGS.37}
\]
There is no index-dependent coefficient choice and
no change of family.  If \(\mathrm{Adm2F}\) is supplied
directly, use its threshold in place of GLG16; the
remaining proof is identical.
% END LGS12

% BEGIN LGS13
\paragraph{LGS13. Full \(\mathrm{LinkG}\) from the remaining common record.}
\label{claim:LGS13}
Under (LGS.2), the same family satisfies
\[
 \bigl(\mathrm{GroupCompatibilityFamily}(\mathcal F)
       \ \wedge\ \mathrm{TopsF}(\mathcal F)\bigr)
               \ \Longrightarrow\ \mathrm{LinkG}(\mathcal F).
                                                               \tag{LGS.38}
\]
One may use \(\mathrm{Adm2F}\) in place of
\(\mathrm{TopsF}\).  There are a constant \(C_G\) and
one threshold \(x_0\), before all later members and
places, such that the error is
\[
           b_o=b^{\delta_G(x)},\qquad
           \delta_G(x)=x^{-\epsilon/2}\longrightarrow0,
           \qquad C_G=4C_{\mathrm{count}}+1.          \tag{LGS.39}
\]
Here \(C_{\mathrm{count}}\ge0\) is exactly the
constant of GLG9 with (LGS.36) and the given \(K\).
No constant or loss term in the literal \(\mathrm{LinkG}\)
is changed.

Apply the complete remaining counting theorem GLG9
to the actual \(X\) and the auxiliary \(Y^0\), whose
full admission is LGS12.  This gives its
\(\mathrm{GainBoundAll}\) and ask comparison for
every pair in the original degree window, in
particular for \(a_L,b_L\).
LGS7 supplies the actual \(\mathrm{MeanKappaBound}\).
LGS8 supplies rational descent, and GLG15 then
identifies all local lengths with the actual
\(L.\mathrm{lenMN}\), including the supported place
with empty served set.  These identifications
continue to hold for \(Y^0\) by LGS11.
The family's original \(\mathrm{GatTies}\),
\(\mathrm{LatTied}\), filling identities, support,
\(\mathrm{NuDef}\), and unserved identity are the
unchanged records on the same member.
GLG10 and GLG12 supply the served-place and depth
conditions.  Finally LGS11 supplies every conjunct
of the local \(\mathrm{admG1}\).
Thus all premises of GLG17 have now been produced.

For precision, its five conclusions here are:
\begin{enumerate}
\item For every \(u,w\) with \(p\ne T_u\), the original
depth inclusion
\[
 \mathrm{LocLe}\bigl(\lfloor\mathrm{ur}_w(\tau(w))\rfloor_+,
                          L.NuLat(u),\mathrm{MramL}(u)\bigr)
 \]
holds, and
\[
 L.\mathrm{lenMN}(u,w)
       \le D.\mathrm{duS}(u)\{\tau(w)+\mathrm{io}(w)\}.
 \]
An empty served set forces all three quantities
\(\tau(w),\mathrm{io}(w),L.\mathrm{lenMN}(u,w)\) to zero.
\item For every prime and served representative,
\(S_v\le2mN_va/p\).
\item At every original signal or inert place with
\(p\le X''\), each gain lies in
\([0,4|TT|a/p]\), and the exact class-sum inequality is
\[
 \begin{aligned}
 &\sum_{j\in\mathrm{lifts}}\ \sum_{t\in TT_w}
       \{\mathrm{vonMangoldtClassSum}(j,t,p)
                         -\mathrm{loG}(L,C_G,b_o,p)\}\\
 &\qquad-m\,\mathrm{flG}(L,C_G,p)
       \le\sum_{j\in\mathrm{lifts}}\mathrm{gain}(j,w)-\tau(w).
 \end{aligned}
 \]
\item The full original \(\mathrm{ServedOrbitsContained}\)
holds.
\item The full original \(\mathrm{MeanKappaBound}\)
holds for \(X\), with both clauses of (LGS.22).
\end{enumerate}
The scalar \(vUnit(G)\) in (LGS.28) cancels by the
exact local length identity in GLG15; the actual
\(\mathrm{GatTies}\) rank identity gives
\(G.d=D.\mathrm{duS}(u)\).
This explains in particular why no coefficient-field
ramification factor is missing in item 1.
The original losses \(\mathrm{loG},\mathrm{flG}\)
and the constant \(4|C_{\mathrm{count}}|+1\) are
the exact GLG11 comparison.  Its absolute value is
unnecessary only because \(C_{\mathrm{count}}\ge0\).

Take the maximum of the finitely many family thresholds
just used.  All five conclusions concern the same
\(x,i,u,w\), lattices, selected tops and degree functions.
Their independence of \(Ch\), checked in LGS11,
turns the auxiliary application into the literal
\(\mathrm{LinkG}(\mathcal F)\).  This proves (LGS.38).

The original class-level implication also follows
with its full choice order.  If a family class
\(\mathcal C\) has
\[
 \mathrm{Closed}(\mathcal C,\mathrm{GroupCompatibilityFamily})
 \quad\hbox{and}\quad
 \mathrm{Closed}(\mathcal C,\mathrm{Adm2F}),
 \]
then it has \(\mathrm{Closed}(\mathcal C,\mathrm{LinkG})\).
First decrease the two supplied positive exponent
cutoffs and \(1/104\) to a common positive cutoff.
For each later fixed \(\epsilon\), take the pointwise
maximum of their two \(\mathrm{TransportedAlt}\)
bounds and
\[
 K_{\mathrm{local}}.\mathrm{pre}=0,\quad
 K_{\mathrm{local}}.t=0,\quad
 K_{\mathrm{local}}.\mathrm{mid}=4129,\quad
 K_{\mathrm{local}}.\mathrm{last}
             =\max\{5600002,C_{\mathrm{orb}}(\epsilon)\}.
                                                               \tag{LGS.40}
\]
The actual \(\mathrm{cMid}\) contains \(K.\mathrm{La}\);
the actual \(\mathrm{cLast}\) contains every other
coefficient constrained in (LGS.2).
Thus \(\mathrm{FamP.Large}\) at this bound implies
(LGS.2) and both supplied class predicates.
Only then is \(\mathcal F\), its height parameter,
and its eventual \(x\)-threshold considered.
This proves the literal source implication
\(\mathrm{linkG\_concrete5}\), with its original
common-family requirement.

The construction of \(\mathrm{GroupCompatibilityFamily}\)
remains outside this theorem.  In particular separate
adelic forms, an abstract local representation, or
unrelated lattice comparisons do not supply its five
simultaneous fields (LGS.24).  The selected-top
assumption is also retained when not provided by
another completed construction.  Consequently this
chapter supplies the local-growth, mean and admission
producers and the exact original installer, while
making no whole-leaf closure claim for the global
local-gain target.
% END LGS13
